\section{System Overview}\label{sec:system}

The proposed workflow links heterogeneous traffic inputs to downstream ETA-oriented evaluation through a single artifact ecosystem. At a high level, the system begins from one of two data-construction paths: a simulation-derived path built on controlled SUMO scenarios, and a real-trajectory path built on imported GPS-like traces aligned to a road network. Although these paths differ substantially in observability, preprocessing, and error sources, both terminate in a common graph-based representation designed for route-aware ETA tasks. This shared representation can then be exercised through a web-based evaluation environment in which a compatible predictor is attached, trip-level outcomes are registered, and statistical analyses are produced.
Figure~\ref{fig:workflow}(\textbf{a}) summarizes this pipeline; panel~(\textbf{b}) sketches the static and dynamic parts of the graph representation.

% Side-by-side: workflow (a) + static/dynamic graph (b); shared height to save vertical space.
\begin{figure}[t]
\centering
\newlength{\pairfigheight}
\setlength{\pairfigheight}{3.75cm}
\begin{minipage}[c]{0.49\linewidth}
\centering
\resizebox{!}{\pairfigheight}{%
\begin{tikzpicture}[
  font=\footnotesize,
  box/.style={
    draw,
    rounded corners=2pt,
    align=center,
    inner sep=3.2pt,
    minimum width=2.55cm,
    line width=0.5pt
  },
  arr/.style={-{Stealth[length=1.6mm]}, thick, shorten >=0.5pt, shorten <=0.5pt}
]
\node[box] (sumo) {\textbf{Sim. path}\\SUMO + TraCI};
\node[box, right=0.28cm of sumo] (traj) {\textbf{Traj. path}\\repair, match};
\coordinate (m) at ($(sumo.south)!0.5!(traj.south)$);
\node[box, below=0.38cm of m] (desk) {\textbf{Desktop}\\\DesktopTool{}};
\node[box, below=0.30cm of desk] (snap) {\textbf{Graph snapshots}\\route-aware, time series};
\node[box, below=0.30cm of snap] (web) {\textbf{Web client}\\interactive ETA evaluation};
\node[box, below=0.30cm of web] (out) {\textbf{Outcomes}\\DB, metrics, reports};
\draw[arr] (sumo.south) to[bend right=10] ([xshift=-0.85cm, yshift=1.2pt]desk.north);
\draw[arr] (traj.south) to[bend left=10] ([xshift=0.85cm, yshift=1.2pt]desk.north);
\draw[arr] (desk) -- (snap);
\draw[arr] (snap) -- (web);
\draw[arr] (web) -- (out);
\end{tikzpicture}%
}\\[-0.05em]
{\scriptsize (\textbf{a}) End-to-end workflow}
\end{minipage}%
\hfill
\begin{minipage}[c]{0.49\linewidth}
\centering
\includegraphics[width=\linewidth,height=\pairfigheight,keepaspectratio]{dynamic_static_nodes_and_edges.png}\\[-0.05em]
{\scriptsize (\textbf{b}) Static vs.\ dynamic graph layers}
\end{minipage}
\caption{(\textbf{a})~End-to-end workflow: simulation and trajectory paths merge in the desktop tool, export route-aware graph time series, and feed the web layer for testing and logged analysis.
(\textbf{b})~Static layer: black nodes A--D are junctions; grey edges are roads. Dynamic layer: green nodes are vehicles; dashed blue edges connect vehicles to other vehicles and to junctions.}
\label{fig:workflow}
\label{fig:graph_schema}
\end{figure}


Dataset construction is operationalized through the \DesktopTool{}, an interactive desktop environment rather than a collection of disconnected scripts. On the simulation side, it supports scenario definition, control of simulation behavior, visual validation of zones, landmarks, and vehicle classes, and export of route-aware graph snapshots with adjustable temporal granularity, including configurable noise such as random untracked trips. On the trajectory side, it supports map download, network preparation, trajectory repair, route alignment, and batch export into the common representation.
Exports are written in a single schema so training code and the \WebApp{} consume the same node/edge feature layout, which reduces mismatch bugs between offline dataset generation and interactive evaluation.

The conceptual center of the system is the common ETA-oriented representation: roads as edges, junctions and vehicles as nodes, with each vehicle associated with a known or inferred route. Traffic state evolves in node and edge features, and graph relations change with vehicle movement. The representation is aimed at vehicle-level ETA but supports other views such as road-level load or area-level aggregates. The \WebApp{} provides the downstream evaluation layer: interactive testing with an attached ETA predictor, outcomes stored in a database, and statistical summaries, plots, and report generation, bridging representation design and measurable use. Compatible predictors can be attached, replaced, and tested in the same environment, including by external users of the open-source release. The cloud-deployed instance discussed later is one demonstrator configuration of the broader workflow.
